\documentclass[lettersize,journal]{IEEEtran}

% ============================================================
% PACKAGES
% ============================================================
\usepackage{amsmath,amssymb,amsfonts,amsthm}
\usepackage{algorithm}
\usepackage{algpseudocode}
\usepackage{graphicx}
\usepackage{cite}
\usepackage{url}
\usepackage{xcolor}
\usepackage{booktabs}
\usepackage{multirow}
\usepackage{bm}
\usepackage[caption=false,font=footnotesize]{subfig}
\usepackage{balance}
\usepackage{tikz}
\usetikzlibrary{arrows.meta,positioning,calc,decorations.pathmorphing,
  shapes.geometric,fit,backgrounds,patterns,fadings,shadows.blur}
\usepackage{pgfplots}
\pgfplotsset{compat=1.18}

\hyphenation{op-tical net-works semi-conduc-tor Kolmo-go-rov}

% ============================================================
% THEOREM ENVIRONMENTS
% ============================================================
\newtheorem{corollary}{Corollary}
\newtheorem{lemma}{Lemma}
\newtheorem{remark}{Remark}
\newtheorem{definition}{Definition}
\newtheorem{proposition}{Proposition}
\newtheorem{assumption}{Assumption}

% ============================================================
% CUSTOM COMMANDS
% ============================================================
\DeclareMathOperator{\erfc}{erfc}
\DeclareMathOperator{\diag}{diag}
\DeclareMathOperator*{\argmin}{arg\,min}
\DeclareMathOperator*{\argmax}{arg\,max}
\newcommand{\vect}[1]{\mathbf{#1}}
\newcommand{\mat}[1]{\mathbf{#1}}
\newcommand{\set}[1]{\mathcal{#1}}
\newcommand{\norm}[1]{\left\|#1\right\|}
\newcommand{\abs}[1]{\left|#1\right|}
\newcommand{\R}{\mathbb{R}}
\newcommand{\E}{\mathbb{E}}

% ============================================================
\usepackage[active,tightpage]{preview}
\setlength\PreviewBorder{0.6pt}
\begin{document}
\begin{preview}
\resizebox{0.8\textwidth}{!}{%
\begin{tikzpicture}[
  node distance=6pt and 8pt,
  cell/.style={circle, draw=black!70, fill=blue!15, minimum size=13pt,
               inner sep=0pt, font=\tiny\bfseries},
  kanbox/.style={rectangle, rounded corners=2pt, draw=#1!70, fill=#1!8,
                 minimum height=28pt, minimum width=14pt,
                 line width=0.6pt},
  kanbox/.default={violet},
  phibox/.style={rectangle, rounded corners=2pt, draw=teal!70, fill=teal!6,
                 minimum height=18pt, minimum width=8pt,
                 line width=0.5pt},
  sumcirc/.style={circle, draw=black!60, fill=yellow!10, minimum size=11pt,
                  inner sep=0pt, font=\scriptsize},
  arr/.style={-{Stealth[length=3.5pt,width=2.5pt]}, line width=0.6pt, black!65},
  darr/.style={-{Stealth[length=3.5pt,width=2.5pt]}, line width=0.8pt,
               blue!60!black},
  lbl/.style={font=\scriptsize, text=black!80},
  tlbl/.style={font=\tiny, text=black!60},
  phase/.style={rectangle, rounded corners=4pt, draw=black!50,
                fill=#1, minimum height=22pt, font=\scriptsize\bfseries,
                align=center, line width=0.6pt},
  phase/.default={gray!8},
]

% ============================================================
% LEFT: Graph input
% ============================================================
\begin{scope}[shift={(0,0)}]
  \node[lbl, font=\small\bfseries, anchor=south] at (1.0,4.1) {Input graph};

  % Small graph
  \coordinate (g1) at (0.2, 3.2);
  \coordinate (g2) at (1.0, 3.6);
  \coordinate (g3) at (0.4, 2.2);
  \coordinate (g4) at (1.7, 2.8);
  \coordinate (g5) at (1.3, 1.7);

  \draw[gray!40, line width=0.7pt] (g1)--(g2) (g1)--(g3) (g2)--(g4)
       (g3)--(g5) (g4)--(g5) (g2)--(g3);

  \node[cell] at (g1) {$1$};
  \node[cell, fill=blue!30] at (g2) {$i$};
  \node[cell] at (g3) {$3$};
  \node[cell, fill=red!15] at (g4) {$j$};
  \node[cell] at (g5) {$5$};

  % Feature vectors
  \node[lbl, anchor=west] at (0.15, 1.05) {$\mathbf{h}_i^{(0)} = \mathbf{c}_i(t)$};
\end{scope}

% ============================================================
% CENTER: RD-GKAN layer (eq. 11)
% ============================================================
% Layer l
\begin{scope}[shift={(3.0,0)}]
  \node[lbl, font=\small\bfseries, anchor=south] at (1.3,4.1)
       {RD-GKAN Layer $l$};

  % Message splines (from neighbors)
  \node[lbl, anchor=south] at (-0.15, 3.55) {from $j$};

  % KAN spline box — message
  \node[kanbox=teal, minimum width=42pt, minimum height=60pt]
       (msgkan) at (1.1, 2.7) {};
  % Fixed MC-derived coupling: no learned nonlinearity on neighbours
  \begin{scope}[shift={(1.10,2.7)}]
    \node[tlbl, teal!70!black] at (0,0.26) {prescribed $w_{ij}$};
    \node[tlbl, teal!70!black] at (0,-0.16) {\emph{fixed, not learned}};
  \end{scope}

  % Aggregation
  \node[sumcirc] (agg) at (2.7, 2.7) {$\Sigma$};
  \node[tlbl, anchor=north] at (agg.south) {over $j$};

  % Self-update spline
  \node[kanbox=violet, minimum width=32pt, minimum height=42pt]
       (selfkan) at (2.7, 1.1) {};
  \node[tlbl, violet!80!black, rotate=90] at (2.35, 1.1)
       {KAN $\bm\Psi^{(l)}$};

  % Spline curve inside self box
  \begin{scope}[shift={(2.6,0.7)}]
    \draw[violet!70!black, line width=0.7pt, smooth]
         plot coordinates {(0,0) (0.06,0.05) (0.13,0.2) (0.22,0.48)
                           (0.3,0.55) (0.38,0.52)};
    \node[tlbl, violet!60!black] at (0.19, 0.72) {$\psi_p^{(l)}$};
  \end{scope}

  % Sum node
  \node[sumcirc] (update) at (4.0, 1.9) {$+$};

  % Arrows
  \draw[arr] (msgkan.east) -- (agg.west);
  \draw[arr] (agg.east) -| (update.north);
  \node[tlbl, font=\tiny, anchor=south] at (3.44,2.78)
       {$-D_\theta[\mathbf{L}_N\mathbf{h}^{(l)}]_i$};
  \draw[arr] (selfkan.east) -| (update.south);
  % identity bypass: the c_i(t) term of eq. (11), routed over the reaction box
  \draw[arr] (2.00,1.1) |- (update.west);
  \node[tlbl, font=\tiny, anchor=south] at (2.95,1.93) {identity};

  % Input/output
  \draw[darr] (-0.7, 2.7) -- (msgkan.west)
       node[midway, above, tlbl]{$\mathbf{h}_j^{(l)}$};
  \draw[darr] (-0.7, 1.1) -- (selfkan.west)
       node[midway, above, tlbl]{$\mathbf{h}_i^{(l)}$};
  \draw[darr] (update.east) -- ++(0.7,0)
       node[midway, above, tlbl]{$\mathbf{h}_i^{(l+1)}$};
  \node[tlbl, font=\tiny, anchor=north] at (1.9,0.44)
       {$\mathbf{h}_i^{(l+1)}=\mathbf{h}_i^{(l)}+\bm\psi_\theta(\mathbf{h}_i^{(l)})-D_\theta[\mathbf{L}_N\mathbf{h}^{(l)}]_i$};
\end{scope}

% ============================================================
% RIGHT: Symbolic recovery
% ============================================================
\begin{scope}[shift={(9.2,0)}]
  \node[lbl, font=\small\bfseries, anchor=south] at (2.1,4.1)
       {Symbolic Recovery};

  % Learned spline plot
  \begin{scope}[shift={(0.0,2.4)}]
    \draw[black!40, ->, line width=0.4pt] (0,0) -- (2.0,0)
         node[right, tlbl]{$c$};
    \draw[black!40, ->, line width=0.4pt] (0,0) -- (0,1.5)
         node[above, tlbl]{$\phi(c)$};
    % Learned spline (sigmoidal, like Hill)
    \draw[blue!70!black, line width=1.2pt, smooth]
         plot coordinates {(0,0.02) (0.2,0.05) (0.4,0.12) (0.6,0.3)
                           (0.8,0.6) (1.0,0.9) (1.2,1.1) (1.4,1.2)
                           (1.6,1.25) (1.8,1.27)};
    \node[tlbl, blue!70!black] at (1.5,0.85) {learned};
    % True Hill function
    \draw[red!60!black, line width=0.8pt, dashed, smooth]
         plot coordinates {(0,0.0) (0.2,0.04) (0.4,0.14) (0.6,0.31)
                           (0.8,0.57) (1.0,0.86) (1.2,1.08) (1.4,1.19)
                           (1.6,1.24) (1.8,1.26)};
    \node[tlbl, red!60!black] at (0.5,0.85) {true};
  \end{scope}

  % Arrow down
  \draw[arr, line width=0.8pt] (1.0, 2.35) -- (1.0, 1.95)
       node[midway, right, tlbl]{LASSO};

  % Bar chart (symbolic projection)
  \begin{scope}[shift={(0.0,0.15)}]
    \draw[black!40, ->, line width=0.4pt] (0,0) -- (3.8,0);
    \draw[black!40, ->, line width=0.4pt] (0,0) -- (0,1.65)
         node[above, tlbl]{$|\hat\alpha_q|$};

    % Bars
    \fill[gray!25]    (0.15,0) rectangle (0.40,0.12);
    \fill[gray!25]    (0.50,0) rectangle (0.75,0.08);
    \fill[gray!25]    (0.85,0) rectangle (1.10,0.05);
    \fill[gray!25]    (1.20,0) rectangle (1.45,0.15);
    \fill[gray!25]    (1.55,0) rectangle (1.80,0.10);
    \fill[red!50!orange, opacity=0.85] (1.90,0) rectangle (2.15,1.40);
    \fill[gray!25]    (2.25,0) rectangle (2.50,0.07);
    \fill[gray!25]    (2.60,0) rectangle (2.85,0.04);

    % Labels
    \node[font=\tiny, rotate=55, anchor=east, text=black!60] at (0.28,-0.07) {$1$};
    \node[font=\tiny, rotate=55, anchor=east, text=black!60] at (0.63,-0.07) {$x$};
    \node[font=\tiny, rotate=55, anchor=east, text=black!60] at (0.98,-0.07) {$x^2$};
    \node[font=\tiny, rotate=55, anchor=east, text=black!60] at (1.33,-0.07) {$x^3$};
    \node[font=\tiny, rotate=55, anchor=east, text=black!60] at (1.68,-0.07) {MM};
    \node[font=\tiny, rotate=55, anchor=east, text=red!70!black]
         at (2.03,-0.07) {\textbf{Hill}};
    \node[font=\tiny, rotate=55, anchor=east, text=black!60] at (2.38,-0.07) {$e^{-x}$};
    \node[font=\tiny, rotate=55, anchor=east, text=black!60] at (2.73,-0.07) {$\ln$};

    % Winner annotation
    \draw[red!60!black, ->, line width=0.5pt] (2.5,1.3) -- (2.18,1.3)
         node[left, xshift=-5pt, tlbl, text=red!70!black, font=\tiny\bfseries]
              {$\hat{\psi}{\propto}\frac{c^2}{K^2{+}c^2}$};
  \end{scope}

  % Epsilon label
  \node[lbl, green!50!black, font=\scriptsize\bfseries] at (2.9, 0.3)
       {$\epsilon_{\mathrm{recon}}$};
\end{scope}

% ============================================================
% Bottom flow arrows connecting the three panels
% ============================================================
\draw[darr, line width=1.0pt] (2.1, 2.5) -- (2.7, 2.5)
     node[midway, above, lbl]{};
\draw[darr, line width=1.0pt] (7.8, 2.2) -- (8.9, 2.8)
     node[midway, above, lbl, rotate=18]{};

% Phase labels at bottom
\node[phase=blue!6, text width=45pt] (phase1) at (1.0, -0.9)
     {Phase 1\\[2pt]\footnotesize Graph\\[1pt]\footnotesize Construction};
\node[phase=teal!6, text width=45pt] (phase2) at (4.5, -0.9)
     {Phase 2\\[2pt]\footnotesize RD-GKAN\\[1pt]\footnotesize Training};
\node[phase=orange!6, text width=55pt] (phase3) at (8.3, -0.9)
     {Phase 3--4\\[2pt]\footnotesize Stability +\\[1pt]\footnotesize Symbolic ID};

\draw[arr, line width=0.7pt] (phase1.east) -- (phase2.west);
\draw[arr, line width=0.7pt] (phase2.east) -- (phase3.west);

\end{tikzpicture}%
}% end resizebox
\end{preview}
\end{document}
